% =====================================================================
% TEMPLATE COTTA — NON compilato da main.tex.
% Per una nuova cotta: copiare in cotte/cottaNN.tex, sostituire i campi
% <...> e i "Da definire", poi aggiungere \clearpage e \input in main.tex.
% Le sezioni "Ricetta di partenza" e "Priming" sono facoltative.
% =====================================================================
\section{Cotta \#N --- <Stile>}
% =====================================================================

\begin{tabularx}{\linewidth}{@{}lX@{}}
\toprule
Data cotta & <gg/mm/aaaa> \\
Stile & <Nome> (BJCP 2021, <codice>, p.~<pagina>) \\
Parametri di stile & OG <a>--<b>; FG <a>--<b>; IBU <a>--<b>; SRM <a>--<b>; ABV <a>--<b>\,\% \\
Ricetta di partenza & <fonte> (<OG>, \SI{<L>}{\liter}) \\
Durata giornata & <x>--<y> \si{\hour} \\
Volume in fermentatore & \SI{<V_f>}{\liter} \\
Efficienza & $\varepsilon = <valore>$ (impianto, cotta \#<n>) \\
Carbonazione obiettivo & $v_t = \SI{<x>}{\vol}$ \\
Imbottigliamento & \textcolor{red}{Da definire} \\
\bottomrule
\end{tabularx}

\subsection*{Obiettivi e risultati}
\begin{tabularx}{\linewidth}{@{}lccX@{}}
\toprule
Grandezza & Obiettivo & Misurato & Note \\
\midrule
$V_{pb}$ [\si{\liter}] & <V_pb> & \textcolor{red}{Da definire} & $V_{fer} + V_{scr}$ \\
$OG$ & <OG> & \textcolor{red}{Da definire} & $GU_{pb} = <GU>$; range dello stile \\
$P_{tot}$ [pt$\cdot$L] & <P> & --- & $GU_{pb} \cdot V_{pb}$ \\
$FG$ & <FG> & \textcolor{red}{Da definire} & $A_{app}$; range dello stile \\
$IBU$ & <IBU> & --- & valore ricetta; range dello stile \\
$ABV$ [\si{\percent}] & <ABV> & \textcolor{red}{Da definire} & prima del priming; $ABV = 131{,}25\,(OG-FG)$ \\
\bottomrule
\end{tabularx}

% ---------------------------------------------------------------------
\subsection*{Ricetta di partenza e modifiche}
\begin{tabularx}{\linewidth}{@{}lXXX@{}}
\toprule
Voce & Partenza (\SI{<L>}{\liter}) & Cotta \#N (\SI{<V_f>}{\liter}) & Motivo \\
\midrule
OG & <OG> & <OG> & \\
Malto & <...> & <...> & \\
Luppolo amaro & <...> & <...> & \\
Luppolo aroma & <...> & <...> & \\
Lievito & <...> & <...> & \\
\bottomrule
\end{tabularx}

% ---------------------------------------------------------------------
\subsection*{Ricetta}

\subsubsection*{Malto}
\begin{tabularx}{\linewidth}{@{}Xrr@{}}
\toprule
Malto & Massa [\si{\kilo\gram}] & \% \\
\midrule
<Malto 1> & <m> & <%> \\
\midrule
Totale & <M> & 100,0 \\
\bottomrule
\end{tabularx}

\noindent Calcolo: $m_i = P_i / (10\, e_i\, \varepsilon)$, con $P_{tot}$ ripartito sui malti. Estratti $e_i$: <...>.

\subsubsection*{Luppolo}
\begin{tabularx}{\linewidth}{@{}Xccccr@{}}
\toprule
Luppolo & Forma & $a$ [\si{\percent}] & Massa [\si{\gram}] & Aggiunta & $IBU_i$ \\
\midrule
<Luppolo> & <pellet> & <a> & <m> & \SI{<t>}{\minute} & <IBU_i> \\
\bottomrule
\end{tabularx}

\noindent Utilizzi: <...>. $C_g$ = <...>.

\subsubsection*{Altri ingredienti}
\begin{tabularx}{\linewidth}{@{}Xcr@{}}
\toprule
Ingrediente & Aggiunta & Quantità [\si{\gram}] \\
\midrule
<...> & \SI{<t>}{\minute} & <q> \\
\bottomrule
\end{tabularx}

\subsubsection*{Lievito}
\begin{tabularx}{\linewidth}{@{}lX@{}}
\toprule
Ceppo & <...> \\
Quantità & \SI{<m>}{\gram} $\;\to\; q_y = \SI{<q>}{\gram\per\liter}$ \\
Attenuazione dichiarata & <...> (usata <...>\,\% per il calcolo di FG) \\
Range temperatura & <...> \\
\bottomrule
\end{tabularx}

\subsubsection*{Acqua}
\begin{tabularx}{\linewidth}{@{}lccc@{}}
\toprule
 & Rocchetta [\si{\liter}] & Lauretana [\si{\liter}] & Totale [\si{\liter}] \\
\midrule
Mash   & <> & <> & <> \\
Sparge & <> & <> & <> \\
\midrule
Totale & <> & <> & <> \\
\bottomrule
\end{tabularx}

% ---------------------------------------------------------------------
\subsection*{Processo}

\subsubsection*{Volumi}
\begin{tabularx}{\linewidth}{@{}lX@{}}
\toprule
Volume totale acqua & $V_{tot} = V_{fer} + V_{scr} + V_{evp} + V_{abs}$ \\
Assorbimento trebbie & $\eta = \SI{1,125}{\liter\per\kilo\gram} \Rightarrow V_{abs}$ \\
Evaporazione & $\sigma = \SI{3}{\liter\per\hour}$ per \SI{60}{\minute} (misurata, cotta \#1) \\
Rapporto di mash $R$ & $\SI{3,5}{\liter\per\kilo\gram} \Rightarrow V_m$ \\
Volume sparge & $V_{sp} = V_{tot} - V_m$ \\
\bottomrule
\end{tabularx}

\subsubsection*{Mash e sparge}
\begin{tabularx}{\linewidth}{@{}Xcc@{}}
\toprule
Step & Temperatura & Durata \\
\midrule
<step> & \SI{<T>}{\celsius} & \SI{<t>}{\minute} \\
\bottomrule
\end{tabularx}

\subsubsection*{Bollitura}
\begin{tabularx}{\linewidth}{@{}lX@{}}
\toprule
Durata & \SI{60}{\minute} \\
Volume post-bollitura & \SI{<V_pb>}{\liter} (obiettivo, prima del raffreddamento) \\
Aggiunte & \SI{0}{\minute}: <...> \\
 & \SI{<t>}{\minute} dopo l'inizio: <...> \\
\bottomrule
\end{tabularx}

\subsubsection*{Fermentazione}
\begin{tabularx}{\linewidth}{@{}lX@{}}
\toprule
Fase 1 & \SI{<T>}{\celsius} per <n> gg \\
Fase 2 & \SI{<T>}{\celsius} per <n> gg \\
Temperatura massima $T_f^{max}$ & \SI{<T>}{\celsius} \\
Durata totale & <n> gg \\
Misurato & \textcolor{red}{Da definire} \\
\bottomrule
\end{tabularx}

\subsubsection*{Priming}
\begin{tabularx}{\linewidth}{@{}lX@{}}
\toprule
Volume CO$_2$ obiettivo $v_t$ & \SI{<x>}{\vol} \\
CO$_2$ residua $v_r(T_f^{max})$ & \SI{<x>}{\vol} (tabella del formulario) \\
Zucchero & $c_s = 3{,}86\,(v_t - v_r) \approx \SI{<x>}{\gram\per\liter}$ \\
$\Delta ABV_p = 0{,}0648\, c_s$ & $\approx <x>$\,\% v/v (ABV finale $\approx$ <x>\,\%) \\
Volume da imbottigliare $V_b$ & \textcolor{red}{Da definire} [\si{\liter}] \\
Massa di zucchero $m_s = c_s V_b$ & \textcolor{red}{Da definire} [\si{\gram}] \\
\bottomrule
\end{tabularx}

% ---------------------------------------------------------------------
\subsection*{Imbottigliamento e prodotto finito}
\begin{tabularx}{\linewidth}{@{}lX@{}}
\toprule
Data & \textcolor{red}{Da definire} \\
Bottiglie & \textcolor{red}{Da definire} \\
ABV finale [\si{\percent}] & \textcolor{red}{Da definire} \\
\bottomrule
\end{tabularx}

% ---------------------------------------------------------------------
\subsection*{Problemi e note per le prossime cotte}
\begin{itemize}
  \item <...>
\end{itemize}
